% Standalone source-record documentation; no research attempts.
\documentclass[11pt]{article}
\usepackage[margin=25mm]{geometry}
\usepackage{fontspec,unicode-math}
\setmainfont{Latin Modern Roman}
\setmathfont{Latin Modern Math}
\usepackage{amsmath,amssymb,mathtools,xurl,hyperref}
\hypersetup{colorlinks=true,urlcolor=blue,linkcolor=blue}
\setcounter{secnumdepth}{0}
\setlength{\parindent}{0pt}
\setlength{\parskip}{5pt}
\setlength{\emergencystretch}{3em}
\newcommand{\sref}[2]{\hyperlink{source:#1:#2}{[S#2]}}
\newcommand{\cataloguescope}[1]{\paragraph{Source record.}#1}
\newcommand{\floor}[1]{\left\lfloor#1\right\rfloor}
\begin{document}
% UnsolvedMath ID 30006218; source catalogue code OWR-14299089-006
\setcounter{section}{30006217}
\section{Duplicate Minimal Convex-Difference Decomposition Problem}\label{30006218}
{\small\sffamily UnsolvedMath ID 30006218 · Computer Science}
\subsection{Definitions and mathematical statement}

\paragraph{Shared notation.}Standard number systems are $\mathbb N,\mathbb Z,\mathbb Q,\mathbb R,\mathbb C$. Specialized notation and index ranges follow the cited source and the definitions below.


For continuous piecewise-linear $f$, each affine region carries a linear formula. A convex-difference representation is $f=g-h$ with convex piecewise-linear $g,h$; minimizing pieces requires counting the regions according to the source\textquotesingle s convention.
\paragraph{Record scope and qualifications.}
The catalogue title is retained as metadata. Named certificates, contracts, bounds or frameworks in that title are not asserted as results of the original source extract.
The supplied registry flags this extraction as unrecoverable. It is displayed as a literal source fragment; the definitions provide context rather than a replacement formal statement.

Source-record limitation: Source flags the statement as unrecoverable: The official report confirms that this overlapping extraction repeats the minimal convex-difference decomposition problem represented canonically by OWR-14299089-005. No separate clean statement is promoted; the original record remains publicly visible and retained for provenance.

Related catalogue records: \href{https://mehmetmars7.github.io/Erdosproblems-llm-hunter/problem.html?type=open_problems&id=30006217}{30006217}. Separate source IDs are preserved for their provenance.
The source text below is rendered literally where necessary; damaged formulas and composite extracts are not silently repaired into a new mathematical claim.
The supplied catalogue marks this record as unrecoverable. Its extract and terminology are documented, but a complete original problem has not been reconstructed.
\paragraph{Statement recorded in the supplied source.}
It remains an open question to construct a minimal decomposition in general dimensions, or to describe a finite procedure to compute it. CWPL function f : Rn \textbackslash{}to R can be written as the difference f = g − h of two convex CPWL functions [2, 4]. Consequently, a natural idea to circumvent the challenges induced by nonconvexity is to use such a decomposition f = g − h and solve the desired problem separately for g and h. However, the crucial question arising from this strategy is: how much more complex are g and h compared to f ? A well-established measure for the complexity of a CPWL function is the number of its linear pieces. Therefore, the main question we study is the following. Problem 1. How to decompose a CPWL function f into a difference f = g − h of two convex CPWL functions with as few pieces as possible? There exist many ways in the literature to obtain such a decomposition, but none of them guarantees minimality or at least a useful bound on the number of pieces of g and h depending on those of f . In fact, no finite procedure is known that guarantees to find a minimal decomposition, with the exception of a recent solution for special cases in dimension 2 [3]. We propose a novel perspective on Problem 1 making use of polyhedral geome- try. Instead of aiming for a globally optimal decomposition, we restrict to Mini-Workshop: Alcoved Polytopes in Physics and Optimization 553 \sref{30006218}{2}
\paragraph{Source status.}
The supplied catalogue labels this source record partially solved; that metadata was checked against the supplied snapshot on 2026-10-11. The label is preserved for provenance and is not an independent verification of present open status or mathematical correctness. The catalogue literature assessment is dated 2026-08-22.
\subsection{Short English statement}
The supplied fragment concerns duplicate Minimal Convex-Difference Decomposition Problem. Its complete mathematical target is not recovered.
\subsection{Sources}
\begin{itemize}
\item[\hypertarget{source:30006218:1}{S1}] Ulam AI and the UnsolvedMath Contributors, UnsolvedMath: A Curated Collection of Open Mathematics Problems, record 30006218 (OWR-14299089-006). Catalogue curation and metadata: CC BY 4.0. The source extract is retained; definitions, English summaries and qualifications are editorial additions.
\url{https://huggingface.co/datasets/ulamai/UnsolvedMath}.
\item[\hypertarget{source:30006218:2}{S2}] Mini-Workshop: Alcoved Polytopes in Physics and Optimization (2025), Oberwolfach Reports, DOI 10.4171/owr/2025/12, pp. 531--558
\url{https://doi.org/10.4171/owr/2025/12}.
\item[\hypertarget{source:30006218:3}{S3}] Brandenburg--Grillo--Hertrich, Decomposition Polyhedra
\url{https://arxiv.org/abs/2410.04907}.
\item Source-specific Oberwolfach Reports licences, including CC BY-SA 4.0 where applicable, are retained; see THIRD\_PARTY\_NOTICES.md.
\end{itemize}
\paragraph{Preparation and review.}
The original source and UnsolvedMath contributors retain attribution for the source extract. Codex (GPT-6; exact serving revision not exposed) assisted with definitions, English summaries, scope notes and formatting. No new human mathematical authorship or independently verified solution is asserted. Permanent source ID 30006218 and code OWR-14299089-006 checked against the complete supplied registry and current repository. Distinct upstream records, including repeated questions and broad or damaged extracts, are retained for complete coverage of the source CS category. A repeated record is not a new independent conjecture. Complete mathematical-equivalence and current-literature review has not been performed. Related canonical identities: 30006217. Preparation and final review follow contribution prompts C and D, with source gaps disclosed. Structural and rendering checks do not establish a new theorem or open status.
\label{end:30006218}
\end{document}
